\subsection{Identity traces over algebraic coefficient fields}
Full spectral retuning is unnecessary for classifying the
base roots of a genuine finite trace. Positivity of their
multiplicities supplies a group-algebra argument. For a
single common spectral radius, identity traces already
give the generic full-centralizer obstruction.

\begin{finding}[\fid{F54}, Part 1: genuine algebraic traces and one radius]
\label{f:algebraicgenuinetrace}
Fix $b>1$ and
\[
 u_n=\sum_{i=1}^s m_i\alpha_i^n,\qquad
 m_i\in\Z_{>0},\quad \alpha_i\in\C^\times,\qquad
 N=\sum_i m_i>0.
\]
Suppose a nonzero polynomial $P(v,Y)\in\C[v,Y]$
of positive $Y$-degree satisfies
\begin{equation}\label{eq:genuinetracealgebraic}
 P(b^n,u_n)=0
\end{equation}
for every sufficiently large integer $n$.
Then there are a common integer $L>0$ and integers
$a_i$ such that
\begin{equation}\label{eq:genuinetracepowers}
 \alpha_i^L=b^{a_i}\qquad\text{for every }i.
\end{equation}
Thus the radii are rational powers of $b$ and the phases
are roots of unity. Neither purity nor an integral
spectral-weight hypothesis is assumed.

Now let $F$ be an invertible finite matrix over an algebraic
extension of $\C(v)$, with a regular invertible value $F_b$.
For every sufficiently large $n$, choose a regular value
of the same family above $b^n$, and suppose only
\begin{equation}\label{eq:genuineidentityretuning}
 \Tr(F_{b^n})=\Tr(F_b^n).
\end{equation}
If all base eigenvalues have one common absolute value
$R>0$, then they are all equal to one number $\lambda$.
If the full centralizer of $F_b$ is semisimple, then
\[
 F_b=\lambda I,\qquad
 \End_{F_b}(\C^N)=M_N(\C).
\]
It therefore cannot be the actual generic algebra
$\C\times\C\times M_2(\C)\times M_2(\C)$.
This conclusion requires no regularity at $v=1$ and
no full characteristic-polynomial comparison.
\end{finding}

\begin{proof}
\par\smallskip\noindent\textbf{A genuine trace in a group algebra.}\quad
Let $\Gamma\subset\C^\times$ be the multiplicative group
generated by $b$ and the finitely many $\alpha_i$.
Its torsion subgroup is finite. Choose a positive integer
$m$ killing that subgroup and put
\[
 \Lambda=\Gamma^m=\{g^m:g\in\Gamma\}.
\]
This group is torsion free: if $g^m$ has finite order,
then $g$ has finite order and hence $g^m=1$.
It is therefore free abelian of finite rank $r$, so
\[
 \C[\Lambda]\simeq
 \C[z_1^{\pm1},\ldots,z_r^{\pm1}]
\]
is an integral domain.
Write $[\lambda]$ for its group-basis monomial and set
\[
 B=[b^m],\qquad U=\sum_i m_i[\alpha_i^m].
\]
The element $B$ is transcendental over $\C$: its powers
are distinct group-basis elements because $b^m$ has
infinite order.

A group-algebra element $W=\sum_\lambda c_\lambda[\lambda]$
gives the finite power sequence
\[
 W(r)=\sum_\lambda c_\lambda\lambda^r.
\]
This correspondence respects sums and products and is
injective even when sequences are compared only for
sufficiently large integers.
Indeed, on distinct nonzero roots
$\lambda_1,\ldots,\lambda_d$, vanishing at $d$ consecutive
indices starting at $r_0$ gives a Vandermonde matrix
times the invertible diagonal matrix
$\operatorname{diag}(\lambda_j^{r_0})$.
Every coefficient must vanish.

Applying this uniqueness to
\eqref{eq:genuinetracealgebraic} on $n=mr$ gives
\begin{equation}\label{eq:genuinegroupring}
 P(B,U)=0\quad\text{in }\C[\Lambda].
\end{equation}
Since $B$ is transcendental, $P(B,Y)$ is nonzero.
Thus $U$ is algebraic over $\C(B)$ in the fraction
field of the Laurent polynomial ring.

\par\smallskip\noindent\textbf{Euler derivations and positive support.}\quad
Choose a lattice basis so that
\[
 B=z_1^d\qquad(d>0).
\]
This is integer linear algebra: the positive gcd of
the coordinates of the nonzero vector for $b^m$ is
$d$, and successive extended-Euclidean coordinate
operations send it to $(d,0,\ldots,0)$ by a
unimodular change of basis.

For $j\ge2$, the Euler derivation
\[
 D_j=z_j\frac{\partial}{\partial z_j}
\]
extends to the fraction field and kills $\C(B)$.
Let $Q(Y)$ be the monic minimal polynomial of $U$
over $\C(B)$. It is separable in characteristic zero,
so $Q'(U)\ne0$.
Applying the derivation to its relation gives
\[
 0=D_jQ(U)=Q'(U)D_jU,\qquad D_jU=0.
\]
After grouping equal Laurent monomials, write
\[
 U=\sum_{a\in S}M_a z_1^{a_1}\cdots z_r^{a_r},
 \qquad M_a\in\Z_{>0}.
\]
The coefficient of this monomial in $D_jU$ is
$M_a a_j$.
Monomial independence and characteristic zero give
$a_j=0$ for every support vector and every $j\ge2$.
All support lies on the first lattice axis.

No original summand disappears during the torsion-killing
or grouping step, because all its multiplicities are
positive. Thus for every $i$ there is an integer $c_i$
with
\[
 [\alpha_i^m]=z_1^{c_i}.
\]
Raising this equality to the $d$th power gives
\[
 [\alpha_i^{md}]=[b^{mc_i}].
\]
Equality of group-basis elements is equality in the
original subgroup of $\C^\times$.
Hence \eqref{eq:genuinetracepowers} follows with
$L=md$ and $a_i=mc_i$.
If the lattice has rank one, the same conclusion holds
without any additional derivation.

For completeness, the relative-algebraic-closure fact
underlying this argument has the same elementary proof.
If $h\in\C(z_1,\ldots,z_r)$ is algebraic over $\C(z_1)$,
separability gives $D_jh=0$ for all $j\ge2$.
The kernel of differentiation in a characteristic-zero
rational function field $M(z)$ is $M$:
for $h=0$ the claim is immediate; otherwise write
$h=p/q$ with coprime polynomials.
The equation $p'q-pq'=0$ implies $p\mid p'$ unless
$p$ is constant; comparison of degrees gives $p'=0$,
and then $q'=0$ as well.
Thus both are constant.
Iterating proves $h\in\C(z_1)$.
Its intersection with the Laurent polynomial ring is
$\C[z_1^{\pm1}]$, also directly seen from monomial
derivatives. The displayed support proof already gives
the precise conclusion needed here.

\par\smallskip\noindent\textbf{The same algebraic trace at one common radius.}\quad
The finitely many entries of the matrix family generate
a finite algebraic extension.
Its trace $f=\Tr F$ is therefore one algebraic element.
Every regular evaluation of that element satisfies
its minimal relation, even if different allowed sheets
are selected at different parameters.
Equation \eqref{eq:genuineidentityretuning} permits the
preceding theorem to be applied to the genuine base
eigenvalues with their algebraic multiplicities.

Choose the common $L$ in \eqref{eq:genuinetracepowers}.
If every base eigenvalue has modulus $R$, then all
the integers $a_i$ are one integer $a$, because
\[
 R^L=b^{a_i}.
\]
Consequently at degrees $n=Lr$,
\[
 \Tr(F_b^{Lr})=N b^{ar}.
\]
The algebraic function
\[
 g(v)=f(v)^L-N^L v^a
\]
therefore vanishes at the selected regular places
above infinitely many distinct values $v=b^{Lr}$.
It must vanish identically.
Indeed, a nonzero algebraic function has a monic
minimal polynomial over $\C(v)$ with nonzero rational
constant term. Outside finitely many coefficient
poles, a zero of the function forces a zero of that
nonzero rational function, which has only finitely
many zeros.
Negative $a$ causes no difficulty in $\C(v)$ or
at the nonzero parameter $b$.
Thus
\begin{equation}\label{eq:commonradiustraceidentity}
 f(v)^L=N^L v^a.
\end{equation}

Evaluate this identity at the original regular place
above $b$, where $f_b=\Tr F_b$.
It gives
\[
 |\Tr F_b|=N b^{a/L}=NR.
\]
This is equality in the triangle inequality for $N$
actual eigenvalues of modulus $R$.
Their genuine positive multiplicities force a common
phase, hence every eigenvalue equals a single $\lambda$.

No diagonalizability has been assumed.
The operator $F_b-\lambda I$ is nilpotent and is
central in the full centralizer.
If that algebra is semisimple, its center is reduced,
so the nilpotent is zero.
The full centralizer is then $M_N(\C)$.
The actual generic algebra of \fid{F19} is semisimple
and nonsimple, giving the stated contradiction.
\end{proof}

\paragraph{The common phase and exact trace-only scope.}
In the unrestricted-radius case the common phase need
not be one.
For example $F(v)=\sqrt v\,I$ may be evaluated on the
negative sheet at $b$, and on the sheets
$(-1)^n b^{n/2}I$ at $b^n$.
The same irreducible algebraic family then obeys
\eqref{eq:genuineidentityretuning} with base eigenvalue
$-\sqrt b$.
Its spectral weight relative to $Q=b^2$ is $1/2$,
so it is outside the integral-weight special case
of \fid{F51}.
The new centralizer conclusion only needs the common
eigenvalue, not removal of this phase.

The group-algebra theorem classifies the base roots
from identity traces alone.
It does not determine the full characteristic
polynomial or eigenvalue functions at other parameters,
nor their collision contact order.
In particular it must not be used to exclude a
multiple-radius model preserving only selected
identity traces.
The stronger contact conclusion of \fid{F52} used
full spectral retuning.

Positive coefficients prevent disappearance of roots
after powering and grouping.
The argument does not recover virtual eigenvalues
that cancel from the visible sequence or a powered
progression.
Arbitrary inserted traces do not automatically
satisfy this positivity hypothesis.
The theorem is a necessary classification, not a
claim that every phase assignment satisfying
\eqref{eq:genuinetracepowers} is realized by one
irreducible algebraic trace function.

The fixed-parameter constructions of \fid{F46},
the faithful tensor families of \fid{F48}, and the
geometric convolution of \fid{F53} remain valid.
Their specified subalgebra actions are different
from equality with a full centralizer under the
additional arithmetic comparison.
No regularity at the collision, integral spectral
weight, or all-conjugate estimate was used for the
common-radius theorem. A pure spectrum at any chosen
embedding supplies its radius hypothesis.

% Separate paragraph for Part III.
\paragraph{\fid{F54}, Part 1, in Grothendieck's language.}
The group algebra $\C[\Lambda]$ is the coordinate
ring of a split torus with character lattice
$\Lambda$. The element $B$ is a character, and
$U$ is the character of a genuine direct sum of
weight lines. Its algebraicity over the one
character forces its support onto that character's
rational line; the lattice index gives the common
power. This is the character/coaction language of
SGA~3, I, 4.7.2--3, with the Laurent and derivative
argument supplied explicitly here.

Under $K_0\operatorname{Rep}(T)=\Z[\Lambda]$, the
element $U$ belongs to the effective cone: its
coefficients are actual weight-line multiplicities.
Virtual classes can cancel them, exactly the step
excluded by the proof.
SGA~6, V, \S\S2--3, and cote~133, scan~18, supply
the previously audited $\lambda/\sigma$ framework
vocabulary \cite{SGA6,Archive}; no $\gamma$-filtration
positivity is inferred.

The algebraic parameter extension is the format
allowed by VIII, \S2.2, p.~7.
The regular matrix and centralizer language uses
EGA~I, Theorem~1.3.7, p.~86, Corollary~1.3.8,
p.~87, and Theorem~1.4.1, pp.~90--92.
Cote~90, scan~17, is only the audited action-framework
anchor \cite{VIII,EGA,SGA3,Archive}.
No manuscript leaf is claimed to contain the new
trace theorem, and no unspecified geometric
Frobenius comparison is being assumed.

